\section{Analytic transfer and repeated binary sweeps}\label{sec:analytic}
\label{ac:part:278}

We now use Theorem~\ref{thm:conditional} with no prescribed paths to bound binary sweeps. Write \(K_\lambda(z)\) for the resulting sweep operator. It is a product of the line averages in \eqref{ac:eq:2}, hence a matrix polynomial in \(z\). All complex parameters below refer to this unconditioned polynomial. Put \(a_0=c_0/(2q)>0\). Since \(c(s)\ge c_0\), \eqref{ac:eq:15} gives
\[
 \|K_\lambda(z)\|_{\op}\le \|K_\lambda(z)\|_{2q}
       \le D_\lambda^{-a_0}\qquad(0\le z\le z_*).
\]

A second estimate holds on a disk of radius greater than one. For an allowed line size \(m\), averaging by \(B_m\) in an irreducible representation is a product of orthogonal projections, one for each independent-switch layer. If the product had norm one, finite dimensionality would give a unit vector attaining that norm. Equality through the successive projections would force this vector to be fixed by every switch group. These groups contain the transpositions along all edges of the connected cube and hence generate \(S_m\). Thus the norm is strictly less than one in every nontrivial irreducible. Let \(\kappa<1\) be the maximum of these norms over the finitely many allowed \(m\) and their nontrivial irreducibles.

Restrict \([\lambda]\) to a line subgroup \(S_m\), and let \(P,Q\) be the averages of \(U_m,B_m\) there. The operator \(P\) is the orthogonal projection onto the line invariants, and \(PQ=QP=P\), because uniform averaging absorbs every group element. Relative to \(\operatorname{ran}P\oplus\ker P\),
\[
 (1-z)P+zQ=I_{\operatorname{ran}P}\oplus z\,Q|_{\ker P},
 \qquad
 \|(1-z)P+zQ\|_{\op}\le\max\{1,|z|\kappa\}.
\]
Indeed \(Q|_{\ker P}\) is a direct sum of the nontrivial line averages, with arbitrary multiplicities, so its norm is at most \(\kappa\). Choose \(r=2\) if \(\kappa=0\), and otherwise choose \(1<r<\kappa^{-1}\). Every line average is then contractive for \(|z|\le r\). Independence makes each stage average the product of its commuting line averages, and the full operator is the ordered product of the stage averages. Submultiplicativity gives \(\|K_\lambda(z)\|_{\op}\le1\) for \(|z|\le r\), uniformly in the grid and \(\lambda\).

To compare this disk bound with the stronger real bound, choose
\(0<u<v<\min\{z_*,1,r-1\}\) and set
\(\Omega=\{z\in\mathbb C:|z|<r\}\setminus[u,v]\). For \(\varepsilon_0>0\), the function
\[
 G(z)=\varepsilon_0\int_u^v\log\frac1{|z-t|}\,dt
\]
is harmonic off the segment. Direct integration of the logarithm gives a continuous extension across the segment, including its endpoints. On \(|z|=r\), the inequality \(|z-t|\ge r-v>1\) gives \(G(z)\le0\). Its values on the closed segment are bounded, so choose \(\varepsilon_0>0\) small enough that \(G\le1\) there. Also \(G(1)>0\), since \(0<1-t<1\) for every \(t\in[u,v]\).

Fix unit vectors \(u_0,v_0\) and put \(f(z)=u_0^*K_\lambda(z)v_0\). This is a scalar polynomial. We may assume \(f\not\equiv0\), since an identically zero coefficient has \(|f(1)|=0\). Then \(\log|f|\) is subharmonic, with value \(-\infty\) at its zeros. Put \(A=a_0\log D_\lambda\ge0\). The function \(H=\log|f|+AG\) is subharmonic on \(\Omega\). The real and disk estimates, together with continuity at the slit, give
\[
 \limsup_{\Omega\ni z\to\xi}H(z)\le
 \begin{cases}
   A G(\xi)\le0,&|\xi|=r,\\
   A\bigl(G(\xi)-1\bigr)\le0,&\xi\in[u,v].
 \end{cases}
\]
The second line includes both sides of the slit and its endpoints; if \(f(\xi)=0\), the logarithm tends to \(-\infty\). The boundary-limsup maximum principle for subharmonic functions~\cite[Chapter I, Section 4.C.1, Theorem (4.14)]{Demailly2012} gives \(H\le0\) on \(\Omega\). Since \(1\in\Omega\), it follows that \(\log|f(1)|\le-A G(1)\). Taking the supremum over the two unit vectors yields
\begin{equation}
 \|K_\lambda(1)\|_{\op}\le D_\lambda^{-g},
 \qquad g=a_0G(1)>0.
 \label{ac:eq:25}
\end{equation}
The constants \(R,q,z_*\) were fixed in Theorem~\ref{thm:conditional}; the choices of \(\kappa,r,u,v,\varepsilon_0\) depend only on them. Thus \(g\) is independent of the grid, its number of coordinates, and \(\lambda\).

It remains to pass from this operator bound to the full permutation law. Let \(\ell=\log_2R\). For \(d\ge\ell\), write \(d=b\ell+t\) with \(0\le t<\ell\), and split the consecutive bits into one block of \(\ell+t\) bits and \(b-1\) blocks of \(\ell\) bits. Their line sizes lie in \([R,R^2]\). At \(z=1\), the binary sweeps on disjoint lines within a block can be interleaved by bit direction, so the grid sweep is exactly \(B_N\) for \(N=2^d\). Its sign coefficient is zero: the sign of any one fair switch is independent of the remaining switches and has mean zero.

For \(\mu_w=B_N^{*w}\), the convolution convention in Section~\ref{sec:prelim} gives \(\widehat\mu_w(\lambda)=K_\lambda(1)^w\). The classical finite-group Fourier argument of Diaconis and Shahshahani~\cite[Section 2, Lemma 2, and the proof of Lemma 14]{DiaconisShahshahani1981} combines Cauchy--Schwarz with matrix Plancherel to give
\[
 \|\mu_w-U_N\|_{\TV}^2
 \le \frac14\sum_{\lambda\ne(N)}D_\lambda\,
                \|K_\lambda(1)^w\|_2^2
 \le \frac14\sum_{\lambda\notin\{(N),(1^N)\}}
                D_\lambda^{\,2-2gw}.
\]
The second inequality uses the sign cancellation and
\(\|K_\lambda(1)^w\|_2^2\le D_\lambda\|K_\lambda(1)\|_{\op}^{2w}\).
A deterministic initial deck translates \(\mu_w\) and leaves its distance from uniform unchanged.

Put \(\beta=2gw-2\) and
\(k(\lambda)=\min(N-\lambda_1,N-\lambda'_1)\). For a fixed \(1\le k\le N/3\), a diagram with \(k(\lambda)=k\) has first row \(N-k\) after transposing if necessary; its remaining \(k\) boxes form a partition of \(k\). There are therefore at most \(2\cdot2^k\) such diagrams, and \eqref{ac:eq:1} gives \(D_\lambda\ge e^{-1}(N/k)^k\). For \(k>N/3\) and sufficiently large \(N\), Section~\ref{sec:prelim} gives \(\log D_\lambda\ge c_{\rm dim}N\), and there are at most \(2^N\) partitions of \(N\). Consequently, for \(\beta>0\) and sufficiently large \(N\),
\[
\begin{aligned}
 \sum_{\substack{\lambda\vdash N\\1\le k(\lambda)\le N/3}}D_\lambda^{-\beta}
 &\le 2e^\beta\sum_{k=1}^{\lfloor N/3\rfloor}
                  \bigl[2(k/N)^\beta\bigr]^k,\\
 \sum_{\substack{\lambda\vdash N\\k(\lambda)>N/3}}D_\lambda^{-\beta}
 &\le \exp\bigl((\log2-\beta c_{\rm dim})N\bigr).
\end{aligned}
\]
For each fixed \(k\), the summand in the first bound tends to zero, and it is at most \((2\cdot3^{-\beta})^k\) throughout its range. Thus the first sum tends to zero when \(\beta>\log2/\log3\), by domination by a summable geometric sequence. The second tends to zero when \(\beta c_{\rm dim}>\log2\). Choose an absolute integer \(w\) so that
\(\beta=2gw-2>\max\{\log2/\log3,\log2/c_{\rm dim}\}\).
The total variation distance then tends to zero as \(d\to\infty\), uniformly over deterministic initial decks.

One binary sweep is exactly \(d\) physical Thorp shuffles, as shown in the introduction. Hence the distance is at most \(1/4\) after \(wd\) physical steps for all sufficiently large \(d\), giving \(t_{\rm mix}(d)=O(d)\). Proposition~\ref{prop:support} gives \(t_{\rm mix}(d)\ge2d-O(1)\), so the mixing time has the optimal order \(\Theta(d)\).
